\section*{Chapitre \ref{ch:ml:training-infrastructure} --- Infrastructure d'entraînement}

\begin{solution}{exo:ml:training-infrastructure:1}
\omterm{def:ml:training-infrastructure:mixed}{bfloat16} $2^{-7}\approx7.8\times10^{-3}$, float16 $2^{-10}\approx9.8\times10^{-4}$, float32 $2^{-23}\approx1.2\times10^{-7}$.
$10^{-30}$ est représentable en \omterm{def:ml:training-infrastructure:mixed}{bfloat16} et en float32 (les plus petits nombres normaux valent environ $10^{-38}$) mais pas en float16, dont
le plus petit nombre positif vaut environ $6\times10^{-8}$~: il passe à zéro par sous-dépassement, raison pour laquelle l'entraînement en float16 demande une mise à l'échelle de la perte.
\end{solution}

\begin{solution}{exo:ml:training-infrastructure:2}
$1\,140 - 20 - 10 + 1 = 1\,111$ fenêtres. Des départs uniformes sur la concaténation produiraient des fenêtres (ou des cibles) qui commencent dans une
séance et finissent dans la suivante, en recollant une clôture à une ouverture~: un saut de cours que le modèle apprendrait à prédire à partir de rien.
\end{solution}

\begin{solution}{exo:ml:training-infrastructure:3}
Avec la projection en mémoire et la boucle, un pas dure $0.74 + 1.17 = 1.91$~ms~; diviser le calcul par deux donne 1.32~ms, soit une accélération de 1.44. Avec le
gather, $0.27 + 1.10 = 1.37$~ms devient 0.82~ms, soit une accélération de 1.67. Plus le chargeur est rapide, plus un \omterm{def:ml:training-infrastructure:accelerator}{accélérateur} plus rapide a de valeur.
\end{solution}

\begin{solution}{exo:ml:training-infrastructure:4}
L'arrondi de chaque pas modifie légèrement la mise à jour~: les deux trajectoires divergent dans l'espace des paramètres et voient les mêmes lots avec
des paramètres différents, de sorte que leurs pertes par pas diffèrent autant que la perte varie d'un lot à l'autre. Les deux convergent vers
des modèles de même qualité~; la comparaison doit porter sur la métrique sur l'ensemble de l'entraînement, et non pas à pas.
\end{solution}

\begin{solution}{exo:ml:training-infrastructure:5}
La perte sur l'union vaut $\frac1{kb}\sum_{j=1}^k\sum_{i\in B_j}\ell_i = \frac1k\sum_jL_j$ où $L_j$ est la moyenne de chaque micro-lot, et
le gradient est linéaire~: la somme des gradients de $L_j/k$ le donne. La \omterm{def:ml:neural-networks-for-noisy-tabular-data:dropout}{normalisation par lot} (batch normalisation) calcule des statistiques par micro-lot, si bien que la
passe avant elle-même diffère et l'équivalence tombe~; il en va de même de tout autre mécanisme qui couple les exemples d'un même lot.
\end{solution}

\begin{solution}{exo:ml:training-infrastructure:6}
Les poids seuls ne déterminent pas le reste de l'entraînement~: les moments de l'optimiseur repartent de zéro, l'échantillonneur rejoue ou saute des données, et les
générateurs aléatoires redémarrent. Dans l'expérience du chapitre, l'entraînement repris s'est écarté d'au plus 0.039 de l'entraînement ininterrompu. Sauvegarder et
restaurer chaque état, et tester la reprise bit à bit.
\end{solution}

\begin{solution}{exo:ml:training-infrastructure:7}
Lignes~: $250\times6.5\times3\,600/0.5 = 11.7$ millions~; binaire $11.7\times10^6\times41\times4 = 1.92$~Go~; texte environ 2.8~Go. À 176~Mo/s
le texte prend environ 16~s par \omterm{def:ml:neural-networks-for-noisy-tabular-data:backprop}{époque}~; les fichiers NumPy, à 1.8~Go/s, environ 1.1~s~; la projection en mémoire ne lit que ce que touchent les lots,
depuis le cache de pages une fois chauffé.
\end{solution}

\begin{solution}{exo:ml:training-infrastructure:8}
Dans un anneau, le tableau est découpé en $p$ morceaux~; une réduction-dispersion fait circuler les morceaux autour de l'anneau $p - 1$ fois, chaque dispositif envoyant
$n/p$ nombres à chaque fois, et une all-gather refait la même chose~: $2n(p-1)/p$ au total, presque $2n$ quel que soit $p$. La communication
par dispositif ne croît pas avec le nombre de dispositifs, mais elle ne décroît pas non plus, alors que le calcul de chaque dispositif décroît comme $1/p$~:
au-delà d'un certain $p$, la communication domine.
\end{solution}

\begin{solution}{pb:ml:training-infrastructure:1}
\textbf{Partie I.}
\begin{enumerate}
\item CSV 2.21~Mo, NumPy 1.50~Mo, projection en mémoire 1.50~Mo.
\item Chargement par \omterm{def:ml:neural-networks-for-noisy-tabular-data:backprop}{époque} 12.5~ms (CSV), 0.83~ms (NumPy), 0.39~ms (projection)~; lot de 0.48 à 0.74~ms en boucle et 0.27~ms avec gather~; calcul de
1.1 à 1.2~ms (mesuré sur l'ordinateur portable de l'auteur).
\item 0.34 (CSV), 0.31 (NumPy), 0.39 (projection, en boucle) et 0.20 (projection, avec gather).
\item La constitution du lot~: un gather vectorisé réduit de moitié la part des données avant que tout changement de stockage ne compte.
\end{enumerate}
\textbf{Partie II.}
\begin{enumerate}[resume]
\item Huit bits d'exposant, sept bits de mantisse~; l'étendue de float32~; epsilon $2^{-7}$.
\item Moyennes des pertes sur les 50 derniers pas 0.0525 (\omterm{def:ml:training-infrastructure:mixed}{bfloat16}) et 0.0527 (float32)~; plus grand écart par pas 0.0132 (22~\%).
\item Son étendue d'exposant est celle de float32~: les petits gradients ne passent pas à zéro par sous-dépassement.
\item La perte, l'état de l'optimiseur et les poids~; les accumulations de nombreux termes (les sommes compensées du livre~4)~; les calculs de risque et de P\&L
en aval du modèle.
\end{enumerate}
\textbf{Partie III.}
\begin{enumerate}[resume]
\item L'accumulation égale un lot plus grand à $1.4\times10^{-7}$ près~; les gradients moyens des fragments égalent le gradient complet à
$1.5\times10^{-8}$ près.
\item Les états du modèle, de l'optimiseur, de l'échantillonneur et des générateurs, ainsi que le pas~; avec les seuls poids, l'entraînement s'est écarté d'au plus 0.039.
\item 693~ko.
\item Redimensionner le \omterm{def:ml:trees-and-boosting:boosting}{taux d'apprentissage} avec la taille du lot, en général avec un échauffement~; et vérifier que le lot plus grand améliore réellement le modèle.
\end{enumerate}
\textbf{Partie IV.}
\begin{enumerate}[resume]
\item \emph{Où part le temps.} Les données représentent 0.34, 0.31 et 0.39 d'un pas d'entraînement avec un stockage CSV, NumPy et projeté en mémoire, et
un lotissement fenêtre par fenêtre, et 0.20 avec la projection en mémoire et un seul gather (mesures sur l'ordinateur portable)~; le plus grand écart de perte entre
float32 et \omterm{def:ml:training-infrastructure:mixed}{bfloat16} sur les 300 pas de l'entraînement est de 0.0132, alors que leurs pertes finales concordent à 0.0002 près.
\item Pas avant le chargeur~: un \omterm{def:ml:training-infrastructure:accelerator}{accélérateur} plus rapide attendrait les données pendant un tiers du temps.
\item Lorsque les paramètres et les activations du modèle ne tiennent pas sur un seul dispositif~: rarement pour les modèles de trading, souvent pour les \omterm{def:ml:large-language-models-in-finance:lm}{modèles de langage}.
\item Fixer les graines, les algorithmes déterministes, les versions des bibliothèques et l'ordre des données~; accepter que des matériels différents arrondissent
différemment et comparer sur les métriques.
\item Les fenêtres à cheval sur une clôture et une ouverture apprennent au modèle un saut qu'il ne peut pas prévoir, et laissent fuir de l'information d'un jour à l'autre.
\item De nombreuses répétitions, des médianes, la charge de la machine consignée, et des rapports plutôt que des durées absolues.
\item La version du code, la version des données, la configuration, les graines, la machine, les \omterm{def:ml:training-infrastructure:ckpt}{points de reprise d'entraînement} (checkpoints) et les métriques (chapitre~25).
\item La part de chaque pas passée à attendre les données.
\end{enumerate}
\end{solution}

\begin{solution}{iq:ml:training-infrastructure:1}
Un format de 16 bits avec l'exposant de float32 et sept bits de mantisse~: moins de précision que float16 mais l'étendue complète, de sorte que les gradients
ne passent pas à zéro par sous-dépassement et qu'aucune mise à l'échelle de la perte n'est nécessaire~; les unités matricielles le font tourner à pleine vitesse.
\iqlookfor{la disposition des bits et l'argument sur l'étendue.}
\end{solution}

\begin{solution}{iq:ml:training-infrastructure:2}
Profiler le pas~: mesurer séparément le chargement des données, le transfert hôte vers dispositif et le calcul~; vérifier le parallélisme du chargeur et la préextraction,
la taille du lot, les points de synchronisation (journalisation, appels à .item()) et les petits noyaux~; corriger le plus gros.
\iqlookfor{mesurer d'abord, puis les coupables habituels.}
\end{solution}

\begin{solution}{iq:ml:training-infrastructure:3}
Chaque dispositif détient une copie du modèle et traite sa part du lot~; après la passe arrière, une \omterm{def:ml:training-infrastructure:dist}{all-reduce} laisse chaque dispositif
avec le gradient moyen, et tous appliquent la même mise à jour.
\iqlookfor{répliques, lots découpés, gradients moyennés.}
\end{solution}

\begin{solution}{iq:ml:training-infrastructure:4}
Des colonnes binaires par instrument et par séance (projetées en mémoire ou découpées en morceaux), un index des frontières de séance, des échantillonneurs qui tirent des fenêtres
à l'intérieur des séances et respectent dans le temps la séparation entraînement/validation, un assemblage des lots en une seule lecture, et une normalisation ajustée sur les
seules données d'entraînement.
\iqlookfor{stockage binaire, échantillonnage conscient des séances et découpages respectant le temps.}
\end{solution}

\begin{solution}{iq:ml:training-infrastructure:5}
Paramètres et tampons du modèle, état de l'optimiseur, état du calendrier du \omterm{def:ml:trees-and-boosting:boosting}{taux d'apprentissage}, position de l'échantillonneur de données et état aléatoire, tous les
états des générateurs aléatoires, le nombre de pas ou d'\omterm{def:ml:neural-networks-for-noisy-tabular-data:backprop}{époques}, ainsi que les versions du code et des données.
\iqlookfor{l'optimiseur et l'échantillonneur, et pas seulement les poids.}
\end{solution}

\begin{solution}{iq:ml:training-infrastructure:6}
Lorsque le modèle dispose de statistiques par lot (\omterm{def:ml:neural-networks-for-noisy-tabular-data:dropout}{normalisation par lot}), lorsque les micro-lots n'ont pas la même taille, lorsque l'ordre de
la sommation en virgule flottante compte pour les résultats bit à bit, et lorsque les dispositifs voient des distributions de données différentes.
\iqlookfor{couches couplées entre les lots, fragments inégaux et ordre de sommation.}
\end{solution}
